Let \(\mathcal L_{\mathcal G}:D(\mathcal I)\to C(\mathcal I)\) map a c\`adl\`ag curve to the piecewise-linear interpolation of its values on \(\mathcal G_{N_0,K_T,K_D}\). For a population panel set
\[
 \widetilde S^0_{p,m}=\mathcal L_{\mathcal G}S^0_{p,m},
 \qquad
 \widetilde s_p=N_0^{-1}\sum_{m\in\mathcal D}\widetilde S^0_{p,m},
 \qquad
 \widetilde Y_m=(\widetilde S^0_{0,m}-\widetilde s_0,
                         \widetilde S^0_{1,m}-\widetilde s_1),
\]
and \(\widetilde\Gamma=N_0^{-1}\sum_{m\in\mathcal D}\widetilde Y_m\otimes\widetilde Y_m\). Apply the selected positive-square-root factor map and pre-period least-squares map described below to \((\widetilde S^0_{p,m})\), with the same fallback rule, and denote the resulting population comparison curve by \(\theta^{\mathrm{pop}}_{r,\mathcal G}\). For \(r=1\), set \(\theta^{\mathrm{pop}}_{1,\mathcal G}=\widetilde s_1\). This comparison curve is a proof object and does not redefine the causal target \(\theta\).

On the grid, compute every full-cell product-limit curve \(\widehat S^{\mathrm{PL}}_{p,m}\), using the convention that an empty product is one, and put \(\widehat S^{\mathrm{KM}}_{p,m}=\mathcal L_{\mathcal G}\widehat S^{\mathrm{PL}}_{p,m}\). For \(r=1\), pool all \(n_D=N_0K_D\) donor-post records before forming the product-limit estimator, denote the resulting c\`adl\`ag curve by \(\widehat S^{\mathrm{PL}}_{D,1}\), and define the continuous estimator by
\[
 \widehat\theta_1=\mathcal L_{\mathcal G}\widehat S^{\mathrm{PL}}_{D,1}.
\]
Writing \(N_{m,j}(u)=\mathbf 1\{X_{1,m,j}\le u,\Delta_{1,m,j}=1\}\), \(Y(u)=\sum_{m,j}\mathbf 1\{X_{1,m,j}\ge u\}\), and \(J(u)=\mathbf 1\{Y(u)>0\}\), define its patient-level multiplier derivative explicitly by
\[
 \widehat Z_1^\xi(t)
 =-\mathcal L_{\mathcal G}\left[
 \widehat S^{\mathrm{PL}}_{D,1}(t-)
 \sum_{m\in\mathcal D}\sum_{j=1}^{K_D}\xi_{1,m,j}
 \int_0^t\frac{J(u)}{Y(u)}\,dN_{m,j}(u)
 \right](t).
\]

For \(r\ge2\), split each donor-period patient set deterministically into odd and even indices. If both folds are nonempty, compute and interpolate the two fold-specific product-limit curves, center each fold across donors to obtain joint-period curves \(\widehat Y_m^{(1)}\) and \(\widehat Y_m^{(2)}\) in \(L^2(\{0,1\}\times\mathcal I)\), and form
\[
 \widehat\Gamma=\frac{1}{2N_0}\sum_{m\in\mathcal D}
 \left\{\widehat Y_m^{(1)}\otimes\widehat Y_m^{(2)}+
 \widehat Y_m^{(2)}\otimes\widehat Y_m^{(1)}\right\}.
\]
Order its eigenvalues \(\widehat\lambda_1\ge\widehat\lambda_2\ge\cdots\), use a deterministic lexicographic sign and basis rule, and set \(\widehat\mu_k=\sqrt{\max\{\widehat\lambda_k,0\}}\). If \(\widehat u_k\) is the corresponding joint-period eigenfunction, put \(\widehat b_{p,k}(t)=\widehat\mu_k\widehat u_{k,p}(t)\) for \(1\le k\le d\), \(\widehat s_p=N_0^{-1}\sum_{m\in\mathcal D}\widehat S^{\mathrm{KM}}_{p,m}\), and \(\widehat H_0=\int_{\mathcal I}\widehat b_0(t)\widehat b_0(t)^\top dt\). On the regular event that both folds are nonempty and
\[
 \widehat\lambda_d\ge\gamma_-^2/4,
 \qquad
 \widehat\lambda_d-\widehat\lambda_{d+1}\ge\gamma_-^2/4,
 \qquad
 \lambda_{\min}(\widehat H_0)\ge\kappa_-/4,
\]
define
\[
 \widehat z_\star=\widehat H_0^{-1}\int_{\mathcal I}\widehat b_0(t)
 \{\widehat S^{\mathrm{KM}}_{0,\star}(t)-\widehat s_0(t)\}\,dt,
 \qquad
 \widehat\theta_r=\widehat s_1+\widehat b_1^\top\widehat z_\star.
\]
This is the coordinate-free selected-positive-square-root prediction map and is invariant to rotations within a tied selected eigenspace. If a fold is empty or any threshold fails, set \(\widehat b_0=\widehat b_1=0\), \(\widehat z_\star=0\), and set
\[
 \widehat\theta_r=\mathcal L_{\mathcal G}\widehat S^{\mathrm{PL}}_{D,1}.
\]
Thus the estimator is continuous and measurable on every dataset in both the regular and fallback branches.

For \(r\ge2\) on the regular event, differentiate the rotation-invariant map from the interpolated cellwise product-limit inputs to its vector of grid predictions. The derivative includes the selected spectral projector, positive square-root, inverse pre-period Gram, target-pre input, and donor means. Apply that derivative to the observed product-limit patient increments and multiply the contribution of patient \((p,m,j)\) by \(\xi_{p,m,j}\). On the fallback event use the explicit pooled derivative \(\widehat Z_1^\xi\) above. Denote the resulting grid process by \(\widehat Z_r^\xi\), and let \(c_r^\xi\) be the conditional \(q_{N_0,K_T,K_D}\)-quantile of \(\max_{g\in\mathcal G}|\widehat Z_r^\xi(g)|\), where
\[
 q_{N_0,K_T,K_D}=1-\alpha+
 \min\left\{\frac{\alpha}{2},
 \frac{\varepsilon_{N_0,K_T,K_D}(\alpha)}{2}\right\}.
\]
The quantile is the left-continuous generalized inverse, so it is defined even when some or all grid-coordinate variances vanish.

Let \(\widetilde\theta_r\) be the pointwise projection of \(\widehat\theta_r\) onto \([0,1]\), with value one imposed at zero. With \(q=\min\{\beta+1,2\}\), \(c_{\mathrm{int}}=c_{\mathrm{int}}(L,\bar\tau,\beta)\) from \(\mathrm{prop:grid\mbox{-}interpolation}\), and
\[
 A_r=c_r^\xi+\mathbf 1\{r\ge2\}\eta_D(K_D)
     +\omega_r(N_0,K_T,K_D;\alpha)
     +c_{\mathrm{int}}h_{N_0,K_T,K_D}^{q},
\]
define
\[
 \widehat B_r^{\mathrm{lo}}(g)=\max\{0,\widetilde\theta_r(g)-A_r\},
 \qquad
 \widehat B_r^{\mathrm{up}}(g)=\min\{1,\widetilde\theta_r(g)+A_r\}.
\]
Interpolate both endpoints linearly and impose their common value one at zero. This defines \(\widehat B_r\) for every dataset. The allowance \(\eta_D\) is the deterministic half-fold risk-exhaustion envelope, \(\omega_r\) is the finite-tail nonlinear and multiplier-transfer buffer, and the off-grid difference between \(S^0\) and \(\widetilde S^0\) is carried only by the final interpolation term.